% !TeX program = lualatex
% !TeX encoding = utf8
% !TeX spellcheck = uk_UA
%
% Лекція 0 «Плазма довкола нас»
% За розд. 0 книги PlasmaBook (PlasmaBook/Nature/Nature.tex). Знімки pictures/ --- копії з PlasmaBook/Nature/Pictures,
% джерела й ліцензії --- у pictures/CREDITS.md.

\documentclass[]{beamer}
\usetheme{PlasmaPhysics}
\usepackage{PlasmaPhysics}
\setlength{\abovedisplayskip}{3pt}
\setlength{\belowdisplayskip}{3pt}
\setbeamercovered{transparent=25}
\graphicspath{{pictures/}}

\title[Лекції з фізики плазми]{\bfseries\huge Плазма довкола нас}
\subtitle{\bfseries Лекції з фізики плазми. Лекція 0}
\author{Пономаренко С. М.}
\date{}

% --- допоміжні команди ---
\newcommand{\pic}[3][0.95\linewidth]{\centering\includegraphics[width=#1,height=#2]{#3}}
\newcommand{\credit}[1]{{\par\tiny\color{gray}#1\par}}

\begin{document}

\usebackgroundtemplate{
	\tikz\node[opacity=0.1]{\includegraphics[width=\paperwidth,height=\paperheight]{background}};%
}

\begin{frame}
	\titlepage
\end{frame}

%% ============================================================
\begin{frame}{План лекції}
	\begin{enumerate}\justifying
		\item Всесвіт як плазма: рекомбінація і реліктове випромінювання.
		\item Природна плазма: Сонце, вітер, магнітосфера, іоносфера, блискавка, міжзоряний газ.
		\item Лабораторна й технічна плазма: розряди, дуга, технології, синтез.
		\item Схожі на вигляд, різні за природою: вогонь, дуга, Сонце й токамак.
		\item Що об'єднує всі ці об'єкти.
	\end{enumerate}
\end{frame}

%% ============================================================
\section{Всесвіт як плазма}
%% ============================================================

\begin{frame}{Плазма~--- стан більшості речовини}
	\begin{block}{}\justifying\small
		За оцінками, понад \alert{99\,\%} видимої речовини у~Всесвіті~--- зірки, туманності, міжзоряний газ~--- перебуває в~стані плазми.
		На Землі плазма трапляється рідше, але значно ближче до нас:
		\begin{itemize}
			\item блискавка і~полярне сяйво,
			\item лампа на стелі, зварювальна дуга,
			\item травлення мікросхем.
		\end{itemize}
	\end{block}
	\begin{alertblock}{Мета лекції}\justifying\small
		Нічого не виводимо: дивимось на знімки й~вчимося бачити фізику, яку з'ясуємо в~наступних розділах.
		Кожен об'єкт~--- точка на діаграмі $n$--$T$ з~розд.~1.
	\end{alertblock}
\end{frame}

\begin{frame}{Рання плазма і рекомбінація}
	\begin{block}{Перші 380 тисяч років}\justifying\small
		Всесвіт заповнювала гаряча \alert{воднево-гелієва плазма}: протони, ядра гелію, електрони й~фотони.
		Фотони безперервно розсіювалися на вільних електронах, тож Всесвіт був \alert{непрозорий}.
	\end{block}
	\begin{block}{Рекомбінація}\justifying\small
		Коли розширення охолодило плазму до $T\sim3\times10^{3}$~К ($\approx0{,}3$~еВ), електрони об'єдналися з~ядрами в~атоми.
		Світло полетіло вільно й~дійшло до нас як \alert{реліктове випромінювання} ($2{,}725$~К).
	\end{block}
\end{frame}

\begin{frame}{Коли Всесвіт рекомбінував}
	\begin{exampleblock}{Оцінка за формулою Саха}\justifying\small
		Сьогодні $n_H\approx2\times10^{-7}~\mathrm{см}^{-3}$. Під час рекомбінації $1+z\approx1100$:
		\[ n\approx2\times10^{-7}\cdot1100^{3}\approx2{,}7\times10^{2}~\mathrm{см}^{-3}. \]
		За Саха $\alpha=1/2$ при $T\approx3{,}7\times10^{3}$~К ($0{,}32$~еВ); при $3\times10^{3}$~К $\alpha\sim5\times10^{-3}$.
		Температура реліктового випромінювання тоді: $2{,}725\cdot1101\approx3\times10^{3}$~К.
	\end{exampleblock}
	\begin{block}{}\justifying\small
		Водень рекомбінує за $T\approx0{,}3$~еВ, а~не $13{,}6$~еВ: на один атом припадає близько $10^{9}$ фотонів,
		і~навіть рідкісні фотони з~<<хвоста>> розподілу іонізують атоми.
	\end{block}
\end{frame}

\begin{frame}{Реліктове випромінювання}
	\begin{columns}[T]
		\begin{column}{0.62\textwidth}
			\pic[\linewidth]{0.55\textheight}{cmb-wmap}
			\credit{WMAP (дев'ятирічні дані), NASA/WMAP Science Team, суспільне надбання}
		\end{column}
		\begin{column}{0.36\textwidth}
			\begin{block}{}\justifying\footnotesize
				Відхилення $\pm2\times10^{-4}$~К від $2{,}725$~К, тобто $\sim10^{-5}$.
			\end{block}
		\end{column}
	\end{columns}
	\begin{block}{}\justifying\small
		Це відбитки \alert{звукових хвиль у~первинній плазмі}: гравітація стискала, тиск випромінювання розпрямляв.
		З~цих згущень виросли галактики.
	\end{block}
\end{frame}

%% ============================================================
\section{Природна плазма}
%% ============================================================

\begin{frame}{Чорні діри та джети}
	\begin{columns}[T]
		\begin{column}{0.48\textwidth}
			\pic[\linewidth]{0.36\textheight}{bh-m87}
			\credit{M87$^{*}$; EHT Collaboration, CC BY 4.0}
		\end{column}
		\begin{column}{0.48\textwidth}
			\pic[\linewidth]{0.36\textheight}{bh-jet}
			\credit{Джет M87; Hubble, NASA/ESA, F.~D.~Macchetto}
		\end{column}
	\end{columns}
	\begin{block}{}\justifying\footnotesize
		Сама діра не є плазмою, але довкола неї~--- \alert{акреційний диск} з~намагніченої плазми ($T$ від мільйонів до мільярдів кельвінів). Кільце на знімку~--- її випромінювання, темна пляма~--- \alert{тінь} діри.
		Маса M87$^{*}$ $\approx6{,}5\times10^{9}M_\odot$, радіус Шварцшильда $\approx130$~а.о.
	\end{block}
	\begin{block}{}\justifying\footnotesize
		\alert{Джети}~--- вузькі потоки плазми зі~швидкостями, близькими до швидкості світла, вздовж осі обертання, до 4~тис.\ світлових років. Їх колімує магнітне поле, закручене разом із~диском.
	\end{block}
\end{frame}

\begin{frame}{Міжзоряний газ. Крабоподібна туманність}
	\begin{columns}[T]
		\begin{column}{0.5\textwidth}
			\pic[\linewidth]{0.55\textheight}{crab-nebula}
			\credit{Hubble; NASA, ESA, J.~Hester, A.~Loll (ASU)}
		\end{column}
		\begin{column}{0.48\textwidth}
			\begin{block}{}\justifying\footnotesize
				Від $\sim10^{-3}$ до $10^{3}$ частинок в~см$^{3}$; $T$ від $10$~К до $10^{6}$~К; газ не врівноважений із~випромінюванням, а~за фронтами ударних хвиль $T_e\ne T_i$.
			\end{block}
			\begin{block}{}\justifying\footnotesize
				Залишок наднової 1054~р., розліт газу на кілька світлових років. Волокна тримає вморожене поле; блакитне світіння~--- синхротронне випромінювання релятивістських електронів.
			\end{block}
		\end{column}
	\end{columns}
\end{frame}

%% ============================================================

\begin{frame}{Сонце і корона}
	\begin{columns}[T]
		\begin{column}{0.5\textwidth}
			\pic[\linewidth]{0.6\textheight}{sun-corona}
			\credit{SDO/AIA, 171~\AA, 14.05.2024, NASA/SDO}
		\end{column}
		\begin{column}{0.48\textwidth}
			\begin{block}{}\justifying\footnotesize
				Ядро: $T\sim1{,}5\times10^{7}$~К, термоядерні реакції.\\
				Фотосфера: $T\approx5800$~К.\\
				\alert{Корона}: $T\sim10^{6}$~К, $n\sim10^{8}~\mathrm{см}^{-3}$ --- гарячіша на два порядки!
			\end{block}
			\begin{alertblock}{}\justifying\footnotesize
				Яскраві петлі~--- плазма, утримана в~магнітних трубках (пучках силових ліній $\vec B$): частинки легко рухаються вздовж поля й~майже не рухаються впоперек.
			\end{alertblock}
		\end{column}
	\end{columns}
\end{frame}

\begin{frame}{Сонячний вітер}
	\begin{block}{}\justifying\small
		Корона не утримується: потік заряджених частинок \alert{сонячний вітер} має біля Землі $300$--$800$~км/с,
		кілька частинок в~одному кубічному сантиметрі, $T\sim10^{5}$~К.
		Вітер несе \alert{вморожене} магнітне поле Сонця; через обертання силові лінії закручуються в~спіраль.
	\end{block}
	\begin{exampleblock}{Скільки маси Сонце втрачає на вітер}\justifying\footnotesize
		$\dot M=4\pi r^2 n m_p v$ при $r=1$~а.о., $n\approx9~\mathrm{см}^{-3}$, $v\approx470$~км/с:
		$\dot M\approx2\times10^{9}$~кг/с, за рік $6\times10^{16}$~кг $\approx3\times10^{-14}$ маси Сонця.
		До Землі вітер летить $\approx4$ доби, корональний викид маси ($\sim10^{3}$~км/с)~--- $\approx1{,}7$ доби.
	\end{exampleblock}
\end{frame}

\begin{frame}{Магнітосфера і полярне сяйво}
	\begin{columns}[T]
		\begin{column}{0.5\textwidth}
			\pic[\linewidth]{0.5\textheight}{aurora}
			\credit{Знімок з борту МКС, NASA}
		\end{column}
		\begin{column}{0.48\textwidth}
			\begin{block}{}\justifying\footnotesize
				Поле Землі ($\sim0{,}5$~Гс) стискає вітер з~боку Сонця й~витягує хвіст \alert{магнітосфери}.
				Перед нею~--- ударна хвиля без зіткнень.
			\end{block}
		\end{column}
	\end{columns}
	\begin{block}{}\justifying\small
		Частинки, спіймані полем, накопичуються в~\alert{радіаційних поясах}; під час бур електрони з~хвоста магнітосфери, прискорені до кількох кеВ, висипаються вздовж силових ліній
		на висоти $80$--$150$~км, і~виникає сяйво (зелена лінія кисню 557,7~нм). Передбачення таких подій~--- \alert{космічна погода}.
	\end{block}
\end{frame}

\begin{frame}{Іоносфера і радіозв'язок}
	\begin{block}{}\justifying\small
		З~висот $\sim60$~км атмосфера іонізована Сонцем: шари D, E, F. Максимум концентрації електронів $10^{5}$--$10^{6}~\mathrm{см}^{-3}$ --- на висотах $250$--$400$~км.
	\end{block}
	\begin{exampleblock}{Хвилі нижче плазмової частоти відбиваються}\justifying\small
		\[ f_p\approx9\sqrt{n}\ \text{кГц},\qquad n=10^{6}~\mathrm{см}^{-3}\;\Rightarrow\;f_p\approx9~\text{МГц}. \]
		Короткі хвилі відбиваються від іоносфери й~приймаються за тисячі кілометрів, ультракороткі проходять крізь неї.
	\end{exampleblock}
\end{frame}

\begin{frame}{Блискавка}
	\begin{columns}[T]
		\begin{column}{0.5\textwidth}
			\pic[\linewidth]{0.5\textheight}{lightning}
			\credit{Техас, NOAA, суспільне надбання}
		\end{column}
		\begin{column}{0.48\textwidth}
			\begin{block}{}\justifying\footnotesize
				Хмара~--- електричний диполь. Поле десь стає достатнім для пробою~$\to$ електронна лавина~$\to$ \alert{лідер} (кроки в~десятки метрів)
				$\to$ зворотний розряд $10^{4}$--$10^{5}$~А.
			\end{block}
			\begin{block}{}\justifying\footnotesize
				Канал: $T$ до $3\times10^{4}$~К; грім~--- ударна хвиля розігрітого газу. Плазма живе соті--десяті частки секунди.
			\end{block}
		\end{column}
	\end{columns}
\end{frame}

%% ============================================================
\section{Лабораторна й технічна плазма}
%% ============================================================

\begin{frame}{Газові розряди}
	\begin{columns}[T]
		\begin{column}{0.48\textwidth}
			\pic[\linewidth]{0.4\textheight}{discharge-tubes}
			\credit{He, Ne, Ar, Kr, Xe; Alchemist-hp, CC BY-SA 2.0 DE}
		\end{column}
		\begin{column}{0.48\textwidth}
			\pic[\linewidth]{0.4\textheight}{plasma-globe}
			\credit{Плазмова куля; Supportstorm, суспільне надбання}
		\end{column}
	\end{columns}
	\begin{block}{}\justifying\footnotesize
		Розряд низького тиску: електрони набирають енергію від поля, а~важкі іони майже ні: $T_e\sim1$~еВ, $T_i\sim300$~К, $n_e\sim10^{10}~\mathrm{см}^{-3}$.
		Колір дає лінійчатий спектр збуджених атомів. У~плазмовій кулі канали~--- \alert{стримери}.
	\end{block}
\end{frame}

\begin{frame}{Дуга і плазмотрони}
	\begin{columns}[T]
		\begin{column}{0.48\textwidth}
			\pic[\linewidth]{0.45\textheight}{welding-arc}
			\credit{Дугове зварювання; W.~M.~Plate~Jr., суспільне надбання}
		\end{column}
		\begin{column}{0.5\textwidth}
			\begin{block}{}\justifying\footnotesize
				Дуга високого тиску: частих зіткнень вистачає, щоб вирівняти $T_e\approx T_i\sim1$~еВ~--- \alert{ізотермічна} плазма, $n_e\sim10^{15}$--$10^{17}~\mathrm{см}^{-3}$.
			\end{block}
			\begin{block}{}\justifying\footnotesize
				\alert{Плазмотрон}: струмінь плазми з~$T=(1$--$3)\times10^{4}$~К, потужність від десятків ват до десятків мегават. Безелектродні розряди дають чисту плазму.
			\end{block}
		\end{column}
	\end{columns}
\end{frame}

\begin{frame}{Технології: різання і напилення}
	\begin{columns}[T]
		\begin{column}{0.48\textwidth}
			\pic[\linewidth]{0.4\textheight}{plasma-cutting}
			\credit{Плазмове різання; D.~Baron Carral, CC0}
		\end{column}
		\begin{column}{0.48\textwidth}
			\pic[\linewidth]{0.4\textheight}{plasma-spray}
			\credit{Напилення; Optimi3e, CC BY-SA 4.0}
		\end{column}
	\end{columns}
	\begin{block}{}\justifying\footnotesize
		Дугу стискає сопло й~потік газу: струмінь $\sim2\times10^{4}$~К розплавляє метал і~здуває розплав (листи до кількох сантиметрів).
		У~струмінь вдувають порошок~--- метал чи кераміка розплавляються й~утворюють покриття.
	\end{block}
\end{frame}

\begin{frame}{Технології: плазмове травлення}
	\begin{columns}[T]
		\begin{column}{0.46\textwidth}
			\pic[\linewidth]{0.5\textheight}{plasma-etching}
			\credit{Поверхня кремнію після травлення, електронний мікроскоп; PolinLinB, CC BY 4.0}
		\end{column}
		\begin{column}{0.52\textwidth}
			\begin{block}{}\justifying\footnotesize
				Газ низького тиску ($1$--$100$~мТор), ВЧ-розряд: $T_e$ --- кілька еВ, $n_e\sim10^{9}$--$10^{11}~\mathrm{см}^{-3}$, сам газ майже холодний.
			\end{block}
			\begin{block}{}\justifying\footnotesize
				Електрони розкладають молекули на \alert{радикали}, поле розганяє \alert{іони} перпендикулярно до пластини: радикали реагують, іони задають напрям~--- вузькі глибокі канавки. Так травлять кремній, алюміній, скло.
			\end{block}
		\end{column}
	\end{columns}
\end{frame}

\begin{frame}{Плазмові двигуни}
	\begin{columns}[T]
		\begin{column}{0.5\textwidth}
			\pic[\linewidth]{0.5\textheight}{hall-thruster}
			\credit{Двигун Холла, 6~кВт; NASA/JPL}
		\end{column}
		\begin{column}{0.48\textwidth}
			\begin{block}{}\justifying\footnotesize
				Іони ксенону розганяє електричне поле в~кільцевому каналі; електрони утримує радіальне $\vec B$ і~вони дрейфують у~схрещених полях $\vec E\times\vec B$.
				Швидкість витікання $10$--$30$~км/с~--- у~кілька разів більша, ніж у~хімічних ракет.
			\end{block}
		\end{column}
	\end{columns}
\end{frame}

\begin{frame}{Гаряча плазма для синтезу}
	\begin{columns}[T]
		\begin{column}{0.5\textwidth}
			\pic[\linewidth]{0.5\textheight}{tokamak-mast}
			\credit{Токамак MAST, Culham; Eye Steel Film, CC BY 2.0}
		\end{column}
		\begin{column}{0.48\textwidth}
			\begin{block}{}\justifying\footnotesize
				Стінка не витримає десятків мільйонів кельвінів, тому плазму \alert{утримує магнітне поле}. У~видимому світлі світиться холодніша периферія, гаряче ядро майже не світить.
			\end{block}
			\begin{block}{}\justifying\footnotesize
				Інерційний шлях: лазерне стиснення мішені. NIF (грудень 2022) --- вихід синтезу вперше перевищив енергію лазера; JET --- рекордні 69~МДж.
			\end{block}
		\end{column}
	\end{columns}
\end{frame}

%% ============================================================
\section{Схожі на вигляд, різні за природою}
%% ============================================================

\begin{frame}{Вогонь і плазма}
	\begin{columns}[T]
		\begin{column}{0.4\textwidth}
			\pic[\linewidth]{0.42\textheight}{campfire}
			\credit{Багаття; T.~Tscharntke, суспільне надбання}
		\end{column}
		\begin{column}{0.58\textwidth}
			\begin{exampleblock}{Чи є полум'я плазмою?}\justifying\scriptsize
				$T\approx2000$~К, $n_e\sim10^{10}~\mathrm{см}^{-3}$, $n\approx3{,}6\times10^{18}~\mathrm{см}^{-3}$: $\alpha\sim3\times10^{-9}$.\\
				$\lambda_{De}\approx3\times10^{-3}$~см $\ll\ell\sim10$~см; $N_D\sim1{,}5\times10^{3}$.\\
				$\omega_p\approx6\times10^{9}~\mathrm{с}^{-1}$, але $\nu\sim10^{11}~\mathrm{с}^{-1}$.
			\end{exampleblock}
		\end{column}
	\end{columns}
	\begin{alertblock}{}\justifying\small
		Критерії колективності виконано, але $\nu\gg\omega_p$: електрон стикається з~нейтралом частіше, ніж робить плазмове коливання. Полум'я~--- гранична слабоіонізована плазма: світіння визначає хімія горіння.
	\end{alertblock}
\end{frame}

\begin{frame}{Розрядна трубка й дуга}
	\begin{block}{}\justifying\small
		Обидві світяться й~є газовим розрядом, але плазми різні: у~трубці електрони рідко стикаються з~важкими частинками, у~дузі~--- часто.
	\end{block}
	\begin{center}\small
		\begin{tabular}{lcc}
			& Розрядна трубка & Дуга \\ \hline
			Тиск & $\sim0{,}1$--$10$~Тор & $\gtrsim1$~атм \\
			$T_e$ & $\sim1$~еВ & $\sim1$~еВ \\
			$T_i$ & $\sim300$~К & $\sim1$~еВ \\
			$n_e$, см$^{-3}$ & $\sim10^{10}$ & $10^{15}$--$10^{17}$ \\
			Ступінь іонізації & $\lesssim10^{-6}$ & відсотки й~більше \\
			Рівновага & $T_e\gg T_i$ & ізотермічна \\
		\end{tabular}
	\end{center}
\end{frame}

\begin{frame}{Сонце й токамак}
	\begin{block}{}\justifying\small
		І~те, й~те~--- гаряча плазма з~термоядерним синтезом. Сонце тримає \alert{гравітація}, токамак~--- \alert{магнітне поле}.
	\end{block}
	\begin{exampleblock}{Скільки потужності дає кубометр Сонця}\justifying\small
		$L\approx3{,}8\times10^{26}$~Вт у~ядрі радіусом $0{,}2R_\odot\approx1{,}4\times10^{8}$~м, $V\approx1{,}1\times10^{25}~\mathrm{м}^3$:
		$L/V\approx30$~Вт/м$^3$, як у~купі компосту. Реактору потрібно близько мегавата на кубометр~--- у~десятки тисяч разів більше.
	\end{exampleblock}
	\begin{block}{}\justifying\small
		Тому на Землі плазма мусить бути гарячішою ($\sim10^{8}$~К проти $1{,}5\times10^{7}$~К) і~горіти на суміші дейтерію з~тритієм.
	\end{block}
\end{frame}

%% ============================================================
\section{Підсумок}
%% ============================================================

\begin{frame}{Що об'єднує всі ці об'єкти}
	\begin{block}{Три властивості}\justifying\small
		\begin{itemize}
			\item \alert{Колективність і~квазінейтральність}: екранування, власні коливання (розд.~1, 2).
			\item \alert{Магнітне поле}: дрейфи, пастки, вморожене поле (розд.~3, 5, 9).
			\item \alert{Рідкісні зіткнення, звичайні хвилі й~нестійкості}: кінетика, теорія нестійкостей (розд.~4, 7, 8).
		\end{itemize}
	\end{block}
\end{frame}

\begin{frame}{Підсумок}
	\begin{block}{}\justifying\footnotesize
		\begin{itemize}
			\item Понад 99\,\% видимої речовини~--- плазма: від первинної плазми (рекомбінація за $T\approx0{,}3$~еВ) до зір і~міжгалактичного газу.
			\item Природна плазма: $n$ від $10^{-5}$ до $10^{26}~\mathrm{см}^{-3}$, розміри від кілометрів до світлових років; структури задає магнітне поле.
			\item Лабораторна: розряд низького тиску нерівноважний ($T_e\gg T_i$), дуга ізотермічна; різання, напилення, травлення, двигуни, синтез.
			\item Полум'я~--- гранична плазма ($\alpha\sim10^{-9}$): $\nu\gg\omega_p$.
			\item Об'єкти розрізняють за $n$, $T_e$, $T_i$, ступенем іонізації та способом утримання.
		\end{itemize}
	\end{block}
\end{frame}

\begin{frame}{Контрольні питання}
	\begin{enumerate}\justifying\small
		\item Чому Всесвіт був непрозорим до рекомбінації й~став прозорим після неї?
		\item Чому водень рекомбінує за $T\sim0{,}3$~еВ, а~не $13{,}6$~еВ?
		\item Чому петлі в~короні повторюють силові лінії?
		\item Чим розрядна трубка відрізняється від дуги за $T_e$, $T_i$, ступенем іонізації?
		\item Чому полум'я~--- слабка плазма? Яка величина це показує?
		\item Чому реактору на Землі потрібна вища температура, ніж у~ядрі Сонця?
		\item Чому іоносфера відбиває короткі хвилі, а~ультракороткі проходять?
		\item Чому в~міжзоряному газі $T_e\ne T_i$?
	\end{enumerate}
\end{frame}

\begin{frame}
	\centering
	\vspace{1.5cm}
	{\bfseries\huge Дякую за увагу!}
\end{frame}

\end{document}
